\documentclass[linenumbers]{aastex701}
\begin{document}

\title{Quantifying the Impact of TESS Light-Curve Provenance on Machine-Learning Variable-Star Classification}
\author{Jonathan Liang}
\affiliation{Adlai E. Stevenson High School, Lincolnshire, IL, USA}
\email{jonathan.liang1223@gmail.com}

\begin{abstract}
Machine-learning analyses of astronomical time series can depend not only on the classifier and extracted features, but also on how the underlying photometry was produced. This study investigates whether the provenance of Transiting Exoplanet Survey Satellite (TESS) light curves is associated with changes in variable-star classification performance and learned feature importance. We constructed a dataset of more than 7,000 variable stars from the AAVSO International Variable Star Index (VSX), spanning nine broad variable-star classes and sampled to provide approximately balanced class representation together with broad coverage in right ascension and declination. A conservative 1\arcsec\ VSX--TESS matching procedure was used to identify suitable TESS observations. When available, light curves from the Science Processing Operations Center (SPOC) or Quick-Look Pipeline (QLP) were used; otherwise, photometry was derived from TESSCut full-frame-image cutouts at the VSX position. An initial Random Forest analysis revealed large differences in classification performance among these provenance groups. To separate provenance-associated effects from differences in the underlying stellar populations, matched-star SPOC--TESSCut and QLP--TESSCut experiments were subsequently constructed in which object identity, class labels, train/test assignments, extracted features, and Random Forest configuration were held fixed. The matched experiments show that classification performance and feature-importance rankings remain dependent on light-curve provenance. These results indicate that heterogeneous TESS light-curve products should not automatically be treated as interchangeable inputs to astronomical machine-learning analyses.
\end{abstract}

% TODO: replace with final AAS Unified Astronomy Thesaurus keywords.
\keywords{Variable stars --- Time series analysis --- Machine learning --- Random forest --- TESS}

\section{Introduction}
\label{sec:introduction}

\subsection{Variable-Star Classification from Photometric Time Series}
Variable stars exhibit changes in observed brightness over characteristic timescales. Their light curves therefore encode information about the physical processes responsible for the variability as well as the observational properties of the source. As modern surveys produce increasingly large collections of photometric time series, automated classification provides a practical way to organize variable sources into scientifically useful categories.

Machine-learning methods provide one approach to this problem. A light curve can be represented by numerical features describing its flux distribution, time-domain behavior, periodicity, and morphology, after which a supervised classifier learns relationships between those features and previously assigned variable-star labels. Random Forest models are useful in this setting because they can represent nonlinear relationships among features and provide measures of feature importance. In this work, feature importance is not interpreted as direct astrophysical causality; rather, it is used to test whether models trained from different light-curve products rely on the same extracted information.

% TODO CITATIONS: variable-star ML; Random Forest; representative TESS classification work.

\subsection{TESS and Heterogeneous Light-Curve Products}
The Transiting Exoplanet Survey Satellite (TESS) provides photometric time-series observations over a large fraction of the sky. TESS observations can be accessed through multiple products and processing pathways. The present study uses three provenance categories: light curves produced by the Science Processing Operations Center (SPOC), light curves produced by the Quick-Look Pipeline (QLP), and light curves derived in this study from full-frame-image cutouts obtained with TESSCut.

These products are related to TESS observations but should not automatically be assumed to be interchangeable representations of an object's variability. Their extraction and processing pathways differ, so measurements associated with the same source can produce different numerical values for variability features. If those features are supplied to a machine-learning classifier, upstream differences can propagate into classification performance and into the model's learned dependence on individual features.

% TODO CITATIONS: TESS mission; SPOC; QLP; TESSCut/MAST.
% Detailed pipeline-processing claims should be added only after verification.

\subsection{The Provenance Problem}
The project initially aimed to construct a Random Forest classifier for variable stars observed by TESS. The input sample was selected from VSX, and each object in the original dataset contributed one usable light curve from SPOC, QLP, or photometry derived from TESSCut. During evaluation of this mixed-provenance dataset, classification performance differed substantially among the three provenance groups. This unexpected result changed the focus of the study from classification alone to a methodological question: \emph{does the provenance of a TESS light curve itself influence downstream machine-learning results?}

The initial observation could not answer this question by itself. Objects with SPOC, QLP, and TESSCut light curves do not form identical stellar samples, and the provenance subsets differ in sample size and class composition. A comparison of their model performance therefore mixes differences in the underlying objects with differences in the light-curve products. The central experimental objective of this work is to separate these factors through matched-star comparisons in which the same objects are analyzed under otherwise identical machine-learning settings.

\subsection{Research Questions and Contributions}
The study is organized around three questions:
\begin{enumerate}
\item Does Random Forest classification performance differ systematically among TESS light-curve provenances?
\item Do models trained from different provenances assign different relative importance to the same set of extracted light-curve features?
\item Do these differences persist when the astrophysical objects, class labels, train/test assignments, feature definitions, and Random Forest configuration are held constant?
\end{enumerate}

To address these questions, we first construct an approximately class-balanced and spatially distributed VSX sample and apply a confidence-first VSX--TESS matching procedure. We then characterize provenance-associated differences in the full dataset and perform two controlled matched-star experiments, SPOC versus TESSCut and QLP versus TESSCut. The purpose is not to introduce a new classification algorithm, but to quantify whether an upstream data-provenance choice changes the behavior of an otherwise fixed machine-learning analysis.

\section{Background and Related Work}
\label{sec:background}

\subsection{Machine Learning for Variable-Star Classification}
Machine-learning classification of variable stars commonly represents each light curve using numerical descriptors and trains a supervised model from objects with known labels. Relevant descriptors can encode the distribution of measured flux, point-to-point variability, periodic structure, and properties of a phase-folded light curve. This feature-based approach is useful for the present study because the same feature definitions can be applied to light curves from different provenances and the resulting models can be compared directly.

Random Forest is the primary classifier used in this work. Logistic Regression was also evaluated as a baseline during model development. The scientific focus, however, is not a comparison of classifier families. Random Forest instead provides a consistent experimental framework in which both predictive performance and learned feature importance can be compared across light-curve provenance. Previous variable-star classification studies establish the broader utility of supervised learning for photometric time series; the present work differs by asking whether changing the upstream TESS light-curve product changes the result obtained from the same downstream analysis.

% TODO CITATIONS: variable-star ML literature; RF; Logistic Regression/scikit-learn; TESS ML.

\subsection{TESS Light-Curve Products}
The three provenance categories enter the analysis through different routes. SPOC and QLP provide pipeline-generated light curves, whereas TESSCut provides image cutouts from which this study derives photometry. Thus, ``provenance'' refers here to the source and processing pathway of the photometric time series supplied to the common feature-extraction and classification pipeline.

\begin{table}[ht]
\centering
\caption{TESS light-curve provenance categories used in this study. Exact processing descriptions will be expanded after verification against the corresponding technical references.}
\label{tab:provenance_overview}
\begin{tabular}{lll}
\hline\hline
Provenance & Input to this study & Light-curve origin \\
\hline
SPOC & Pipeline light curve & SPOC product \\
QLP & Pipeline light curve & QLP product \\
TESSCut & FFI cutout & Photometry derived in this study \\
\hline
\end{tabular}
\end{table}

This difference in processing pathway provides a plausible route by which provenance can influence downstream features. The present paper therefore treats provenance as an experimental variable rather than assuming that all TESS light-curve products are interchangeable.

% TODO CITATIONS: SPOC; QLP; TESSCut.

\subsection{VSX--TESS Cross-Matching and Dataset Construction}
The AAVSO International Variable Star Index provides the source catalog and class labels used to construct the supervised-learning sample. The present study begins from VSX rather than from a pre-existing TESS light-curve catalog. This direction of construction allows the target sample to be selected intentionally across variable-star families, subtypes, and sky position before TESS availability and provenance are determined.

The selected VSX coordinates are associated with TESS sources using the conservative positional-matching procedure described in Section~\ref{sec:dataset}. The method uses a 1\arcsec\ search radius for TIC association and prioritizes a high-confidence positional match. If the matched source has a suitable SPOC or QLP light curve, that product is used; otherwise, TESSCut is used to obtain an image cutout at the original VSX coordinates and photometry is derived from that cutout. Proper-motion checks using SIMBAD are used as an independent validation of the coordinate-based matching strategy.

% TODO CITATIONS: VSX; TIC; SIMBAD; Fetherolf et al.

\subsection{Data Provenance as a Methodological Variable}
A machine-learning model operates on measurements that have already passed through an observational and processing chain. Differences introduced before feature extraction can therefore propagate into the statistical representation seen by the classifier. If provenance changes predictive performance or feature-importance structure, combining heterogeneous products without accounting for their origin can introduce a methodological dependence into the final analysis.

This study tests that dependence empirically. The initial full-dataset analysis establishes whether model behavior is associated with provenance, while the later matched-star experiments test whether the differences persist after major stellar-population confounds are controlled.

% TODO: targeted literature review for provenance/pipeline dependence in astronomical ML.

\section{Data and Dataset Construction}
\label{sec:dataset}

The dataset was designed for supervised classification while also supporting a later investigation of TESS light-curve provenance. Construction began with labeled variable stars in VSX rather than with objects selected from a particular TESS pipeline. The procedure consisted of four stages: defining a broad variable-star class taxonomy, sampling VSX objects across class and sky position, associating those objects with TESS observations using a conservative matching procedure, and validating the coordinate associations before downstream feature extraction.

\subsection{VSX Source Catalog and Class Taxonomy}
\label{subsec:vsx_taxonomy}
The source population was drawn from VSX. The detailed VSX variability types were consolidated into nine broad families used as the supervised-learning labels: \texttt{CEPHEID}, \texttt{CV}, \texttt{DSCT\_SXPHE}, \texttt{ECLIPSING}, \texttt{ELLIPSOIDAL\_ROT}, \texttt{LONG\_PERIOD}, \texttt{RRLYR}, \texttt{XRAY}, and \texttt{YSO}.

The sampling procedure imposed class-dependent target caps rather than assuming that every family contained the same number of suitable objects. The working caps were 1000 objects for \texttt{CEPHEID}, \texttt{DSCT\_SXPHE}, \texttt{ECLIPSING}, \texttt{ELLIPSOIDAL\_ROT}, \texttt{LONG\_PERIOD}, \texttt{RRLYR}, and \texttt{YSO}; 603 for \texttt{CV}; and 62 for \texttt{XRAY}. These unequal caps reflect the available source populations and mean that the dataset is more accurately described as approximately class-balanced rather than exactly balanced.

\begin{table}[ht]
\centering
\caption{Variable-star families and working VSX sampling caps. Final post-processing class counts should be inserted from the canonical dataset before manuscript submission.}
\label{tab:class_taxonomy}
\begin{tabular}{lc}
\hline\hline
Class family & Sampling cap \\
\hline
\texttt{CEPHEID} & 1000 \\
\texttt{CV} & 603 \\
\texttt{DSCT\_SXPHE} & 1000 \\
\texttt{ECLIPSING} & 1000 \\
\texttt{ELLIPSOIDAL\_ROT} & 1000 \\
\texttt{LONG\_PERIOD} & 1000 \\
\texttt{RRLYR} & 1000 \\
\texttt{XRAY} & 62 \\
\texttt{YSO} & 1000 \\
\hline
\end{tabular}
\end{table}

% TODO: add exact VSX subtype-to-family mapping from the canonical loader/notebook.
% TODO CITATION: AAVSO VSX.

\subsection{Stratified VSX Sampling}
\label{subsec:vsx_sampling}

A simple random draw from VSX could produce a training sample dominated by the most numerous categories or concentrated in particular regions of the sky. The selection procedure was therefore designed to distribute objects across variable-star family, subtype, right ascension, and declination. The goal was not to reproduce the natural frequency of each VSX class, but to construct a supervised-learning dataset in which multiple variability families were represented sufficiently well for model training and provenance-stratified evaluation.

Within each broad family, the selection procedure considered available VSX subtypes and sky coordinates when selecting objects up to the relevant class cap. This provided broad right-ascension and declination coverage while reducing the chance that a small number of highly populated subtypes would dominate an entire family. Because some classes contained fewer available objects than the nominal 1000-object target, exact equality among classes was neither possible nor required.

This sampling choice is important for interpreting later performance metrics. Overall accuracy is sensitive to class prevalence, and the final provenance subsets are not identical in class composition. Later analyses therefore report balanced accuracy in addition to ordinary accuracy and, more importantly, use matched-star experiments to remove differences in object identity and class composition when testing the provenance hypothesis.

% TODO: replace/expand with pseudocode-level details from the canonical
% VSXCategoryLoader/data-pipeline implementation before manuscript lock.

\subsection{Confidence-First VSX--TESS Matching}
\label{subsec:crossmatch}

For each selected VSX object, the pipeline attempted to associate the catalog position with a TESS Input Catalog (TIC) source using a 1\arcsec\ search radius. The matching strategy was deliberately conservative: the closest TIC source within the allowed radius was treated as the candidate association. The 1\arcsec\ threshold was chosen to prioritize positional confidence rather than maximizing the number of possible catalog associations.

After a TIC association was established, the pipeline determined whether a suitable pipeline light curve was available. When the selected TIC had an available SPOC or QLP product, that pipeline light curve was adopted for the original mixed-provenance dataset. When a suitable pipeline light curve was not available, the procedure fell back to TESSCut. In the fallback case, the TESSCut extraction was centered on the original VSX sky position rather than requiring a lower-confidence TIC association. Photometry was subsequently derived from the returned cutout as described in the light-curve processing section of the manuscript.

Conceptually, the procedure can be summarized as
\begin{equation}
\mathrm{VSX\ position}
\rightarrow
\begin{array}{ll}
\mathrm{matched\ TIC\ with\ suitable\ SPOC/QLP} & \rightarrow \mathrm{pipeline\ light\ curve},\\
\mathrm{otherwise} & \rightarrow \mathrm{TESSCut\ at\ VSX\ position}.
\end{array}
\end{equation}

This design served two purposes. First, it avoided broadening the positional criterion merely to obtain a pipeline product, thereby keeping the catalog association conservative. Second, the TESSCut fallback allowed objects without a suitable SPOC or QLP product to remain in the dataset. The resulting sample therefore contains heterogeneous light-curve provenance by construction, which later made it possible to identify the strong provenance-associated differences that motivated the controlled experiments.

Each object in the original classification dataset contributed one light curve and one provenance label. Thus, the initial dataset did not contain multiple representations of the same star for the primary mixed-provenance model. Multiple representations were generated later only for the matched-star provenance experiments, where TESSCut light curves were deliberately produced for the same objects already represented by SPOC or QLP.

% TODO: verify exact SPOC-versus-QLP priority rule and availability logic from
% the canonical converter/downloader before finalizing this paragraph.
% TODO CITATIONS: TIC; TESSCut; SPOC; QLP.

\subsection{Cross-Match Validation}
\label{subsec:match_validation}

Because an incorrect VSX--TIC association could propagate through every later stage of the analysis, the positional matching procedure was evaluated independently of the machine-learning results. The validation considered positional separation, duplicate or ambiguous catalog associations, and the possible effect of stellar proper motion between the catalog coordinate epoch and the TESS observation epoch.

\subsubsection{Positional Separation and Duplicate Associations}
\label{subsubsec:positional_validation}

For TIC-associated objects, the angular separation between the VSX position and selected TIC candidate provides the most direct diagnostic of the positional match. By construction, accepted associations lie within 1\arcsec, and the closest candidate satisfying this criterion is selected. The final manuscript should report the distribution of accepted separations rather than relying only on the maximum allowed radius. Useful summary quantities include the median separation, upper percentiles, number of successful TIC associations, and the number of cases in which the same TIC identifier is associated with more than one selected VSX object.

Duplicate-TIC checks are particularly important because a very small search radius does not by itself guarantee a one-to-one mapping if catalog entries are crowded or duplicated. These diagnostics should therefore be reported alongside the separation distribution to establish the practical reliability of the matching procedure.

% TODO RESULTS: insert exact match-success count, separation median/percentiles,
% and duplicate-TIC statistics from the canonical matching analysis.
% No numerical values are invented in this draft.

\subsubsection{Proper-Motion Validation}
\label{subsubsec:proper_motion}

A second validation considered whether proper motion could make a fixed-radius coordinate match unreliable. The working coordinate treatment assumes the catalog sky position at its reference epoch, whereas TESS observations occur at later epochs. For a star with proper-motion components $\mu_{\alpha *}$ and $\mu_{\delta}$, the approximate displacement over a time interval $\Delta t$ can be written as
\begin{equation}
\Delta\theta \simeq
\Delta t\sqrt{\mu_{\alpha *}^{2}+\mu_{\delta}^{2}},
\label{eq:proper_motion_displacement}
\end{equation}
with consistent angular and temporal units.

SIMBAD proper-motion information was retrieved as an independent check on this issue. TESS sectors were mapped to their observation epochs, allowing the expected displacement between the catalog reference epoch and the relevant TESS observation to be estimated for objects with available proper-motion measurements. The resulting displacement distribution can then be compared with the 1\arcsec\ matching radius.

This check is intended to test the robustness of the coordinate-based association rather than to redefine the matching algorithm on an object-by-object basis. In the final manuscript, the proper-motion analysis should report the number of objects with usable SIMBAD measurements, the typical and upper-percentile predicted displacement, and the number of objects---if any---for which the expected displacement approaches or exceeds the adopted matching radius.

% TODO RESULTS: insert exact SIMBAD coverage and displacement statistics.
% TODO: verify the exact reference epoch used by the VSX-coordinate pipeline.
% TODO CITATION: SIMBAD.

\subsection{Final Dataset}
\label{subsec:final_dataset}

After catalog selection, TESS association, light-curve retrieval or extraction, and quality filtering, the working dataset used for the provenance analysis contained 7,281 variable stars. The provenance distribution reported in the current analysis is 509 SPOC light curves, 1,353 QLP light curves, and 5,419 light curves derived from TESSCut. These counts sum to the full working sample and demonstrate that the provenance groups are substantially unequal in size.

\begin{table}[ht]
\centering
\caption{Working provenance composition of the full dataset. These values should be checked once more against the canonical final data products before manuscript submission.}
\label{tab:provenance_counts}
\begin{tabular}{lr}
\hline\hline
Provenance & Number of stars \\
\hline
SPOC & 509 \\
QLP & 1,353 \\
TESSCut & 5,419 \\
\hline
Total & 7,281 \\
\hline
\end{tabular}
\end{table}

The unequal provenance counts are an important feature of the experimental design rather than a nuisance to be ignored. SPOC, QLP, and TESSCut availability is not distributed identically across the selected VSX population, so the three provenance subsets differ in both sample size and class composition. Consequently, comparisons made directly among these full subsets can reveal a provenance-associated pattern but cannot isolate provenance from differences in the underlying stellar populations.

For this reason, the full 7,281-object dataset serves two roles. It provides the original supervised-learning sample in which the provenance effect was discovered, and it provides the source population from which the later controlled experiments were motivated. The causal interpretation of the provenance comparison does not rest on these unequal full subsets alone. Instead, the central test uses matched objects represented by different light-curve products while holding the labels, train/test assignments, extracted features, and Random Forest configuration fixed.

% TODO: add a class-by-provenance table from the final canonical dataset.
% TODO: verify whether "after quality filtering" is the exact stage represented
% by the 7,281 count; adjust wording if the canonical notebook defines it differently.

\subsection{Dataset-Construction Summary}
\label{subsec:dataset_summary}

The construction procedure was designed to separate two decisions that are often coupled in astronomical machine-learning datasets: which astrophysical objects should be included and which photometric product should represent them. Objects were first selected from VSX according to the variable-star taxonomy and spatial sampling strategy. TESS provenance was then determined through the conservative matching and fallback procedure. This ordering allowed the supervised-learning population to be defined independently of any single TESS pipeline.

At the same time, the resulting full dataset remained heterogeneous in provenance. That heterogeneity became scientifically important when the initial classifier showed markedly different performance across SPOC, QLP, and TESSCut objects. Sections that follow describe the common light-curve processing and feature extraction applied before classification and then distinguish the initial unmatched provenance analysis from the matched-star experiments designed to control the underlying stellar population.

% -----------------------------
% End of requested draft.
% Next planned section:
% \section{Light-Curve Processing and Feature Extraction}
% -----------------------------


% Bibliography
\bibliographystyle{aasjournalv7}
\bibliography{references_fixed}

\end{document}
